\section{Introduction}

Major histocompatibility complex class I (MHC-I) molecules display intracellularly derived peptides at the cell surface for surveillance by CD8\(^{+}\) T cells. The resulting immunopeptidome is shaped by substantially more than peptide--MHC binding: protein abundance and turnover, proteolytic degradation, antigen-processing pathways, cellular composition, tissue environment, and experimental detectability can all affect which peptides are observed. Multi-tissue immunopeptidomic studies have consequently revealed both widely shared ligands and pronounced tissue-associated differences \cite{marcu2021atlas,schuster2018mouse}. Nevertheless, most computational models are designed to predict peptide--MHC binding or presentation in a general setting and therefore do not directly characterize how presentation evidence varies across tissues under a fixed MHC context.

Existing computational approaches can be broadly grouped into three categories. Binding predictors estimate peptide--MHC compatibility or affinity; eluted-ligand models learn general presentation from mass-spectrometry-derived ligands; and expression-aware or sample-context models augment presentation prediction with molecular measurements or sample information \cite{reynisson2020netmhcpan,odonnell2020mhcflurry,albert2023bigmhc}. These formulations have advanced antigen prioritization, but they do not establish an explicit benchmark in which the prediction target is a tissue--MHC task and positive evidence is contrasted with peptides matched by both MHC restriction and parent protein. As a result, they do not directly answer whether two peptides originating from the same protein and observed under the same MHC restriction show different presentation-evidence preferences in a specified tissue.

We formulate this question as \emph{tissue--MHC-specific presentation preference prediction}. Let \(p\) denote a peptide and let \(\tau=(t,m)\) denote a target task defined by tissue \(t\) and MHC restriction \(m\). A model produces a real-valued score
\[
s_{\theta}(p,\tau)\in\mathbb{R},
\]
which is used to rank presentation-evidence preference within task \(\tau\). This is a closed-task setting: the proposed model uses task-specific prediction heads and is primarily intended for tissue--MHC tasks observed during training. It should therefore not be interpreted as an unseen-tissue or unseen-MHC predictor.

To construct supervision aligned with this objective, each positive peptide is paired with a pseudo-negative peptide under the constraints
\[
m^{+}=m^{-}, \qquad u^{+}=u^{-},
\]
where \(m\) is the MHC restriction and \(u\) is the parent UniProt protein. The pseudo-negative has positive evidence in another tissue but is not reported for the target tissue--MHC task. Matching MHC restriction removes variation in the nominal presenting molecule, while matching the parent protein partially controls protein identity, tissue-expression differences, and study bias. The resulting comparison forces the model to distinguish peptides from the same source protein according to task-conditioned evidence. Importantly, absence of a report is not proof of non-presentation, and this paired design cannot fully eliminate expression, study, batch, or mass-spectrometry detection biases.

On this basis, we introduce the TissuePMHC benchmark, a paired human and mouse resource for tissue-conditioned MHC-I presentation preference, together with a structured human model. To avoid conflating the resource and the architecture, we use \emph{TissuePMHC benchmark} for the task and datasets and \emph{TissuePMHC human model} for the multi-kernel dual-branch predictor. The human model combines a position-preserving multi-kernel peptide encoder with complementary global and HLA-specific branches, auxiliary supervision, fixed task-wise rank fusion, and multi-seed averaging. The human and mouse benchmarks are evaluated under both standard pair-disjoint and connected-component peptide-disjoint protocols. The former measures performance in the primary closed-task benchmark, whereas the latter tests seen-task, unseen-peptide behavior while preventing any peptide entity from crossing folds.

Our contributions are fivefold. First, we formalize tissue--MHC-specific presentation preference and construct MHC- and parent-protein-matched human and mouse benchmarks. Second, the TissuePMHC human model aligns global sharing, HLA-structured sharing, auxiliary learning, and multi-scale peptide encoding with the human task structure. Third, matched standard and peptide-disjoint evaluations quantify task learnability and the unseen-peptide generalization gap in both species. Fourth, architecture controls trained on identical peptide-component folds distinguish strict task learnability, retained component benefits, and multi-seed ensemble effects. Fifth, capacity-matched MHC-only models and frozen MHCflurry/NetMHCpan controls quantify how much of the benchmark signal can be recovered without tissue identity.
